\section{Revision and Scope}
We thank the referee for distinguishing the genuine progress in the R7 evidence from the missing integrated paper and the unresolved role of the action oracle. This revision retains the title \emph{Neural Bellman Operators}, the original preference-adjustment economy and utility normalization, recursive utility, temporal selves, dynamic competition, viscosity selection, and coupled nonlinear capital. We have revised the paper itself, rather than replacing the paper by a protocol or treating a successful workflow as an economic result.

The revision is based on the advisory report of 3 October 2026 at commit \texttt{c63f6817f498e1cfe9a653139f7a25bf44a4418a}. The report was repository-owner commissioned and is not an Econometrica editorial decision. R7 source and evidence are pinned separately. The root article and supplement now incorporate the actual R7 results and the executed R8 studies. The preceding sources are archived by exact Git blob identity. A generated-table manifest and execution manifest identify the evidence used in every new numerical table.

The new work has four parts: continuous-action actor training without exhaustive maximizing labels; an operational action-box upper envelope that remains valid at a finite verification budget; matched additional direct, classical, gated and projection comparisons; and complete integration of the finite-policy, strategic, coverage, arithmetic-replay, and failure records. We distinguish an achieved numerical result from a condition needed to transfer it to a continuous economy.

\section{Blocking Comments}
\subsection{B1: Integration of the manuscript}
The referee is correct that the R7 evidence branch did not change the root article and supplement. Both are now revised. The abstract, introduction, algorithmic discussion, preference application, game application, coverage analysis, conclusion, and replication appendix identify the computations actually performed. New tables are generated from result files, not entered as narrative targets. The supplement contains proofs, all method-paired comparisons, missing-pure-stage locations, frozen-policy refinements, and the earlier R6 numerical discussions explicitly marked as historical. All economic themes and earlier evidence remain available.

The candidate is self-contained: it includes the article, supplement, this response, numerical source, raw arrays and network snapshots, logs, input and output hashes, table checks, and compilation records. The later unexecuted R7 development protocol is not represented as evidence. The final candidate branch is distinct from its immutable computational source.

\subsection{B2: Exhaustive action labels and the actor's role}
We now state explicitly that the R7 finite NDU and game actors were trained using exhaustively computed payoff vectors and exact finite maximizing or minimum-defect labels. Their successful raw policies establish neural representation quality on the finite models; they do not establish avoided maximization.

We additionally execute a different training route in the original preference economy. Consumption, portfolio exposure, and preference adjustment are continuous bounded outputs. The actor is updated through its own positive-weight continuation payoff, without constructing an all-action payoff matrix or receiving an exact maximizing label. Reference and global action-cover calculations begin only after the complete policy has been frozen. The code records zero exhaustive training backups and zero exact argmax training labels.

The initial continuous actor has an economically material local failure, retained in the record. A disclosed continuous local-improvement step repairs much of that loss. We therefore report the raw actor, hybrid actor-plus-search, direct continuous optimization, and actor-free search separately. Local improvement changes nearly all principal active nodes, so the hybrid is not described as an actor rescued at only a few exceptional states. This experiment removes the exhaustive \emph{training} oracle, but it does not by itself establish an actor-only advantage. The complete verification oracle remains independent and fully costed.

\subsection{B3: Controlled method comparisons}
The unfavorable R7 NDU comparison is now printed in the main text: at the recorded budget, direct neural fitting is faster and more accurate on average, and classical backward induction is far cheaper on the small two-state grid. We neither omit those methods nor select a favorable seed.

R8 adds eight-restart direct continuous-control optimization and an actor-free local-search ablation with the same continuation class and local improvement routine. It also executes gated-network finite Bellman regression, tensor Chebyshev projection at three degrees, actor-free neural dynamic-game continuation fitting, and classical dynamic-game selection. Every game comparator receives an independent full dynamic best response. Gated regression is explicitly a DGM-style \emph{architecture} comparison on finite Bellman targets, not a claim to have run the mesh-free differential DGM algorithm. A deep-BSDE implementation is not falsely included in the measured set.

Nominal update matching is distinguished from parameter count, actual wall time, closure count, and achieved economic error. The complete-action table reports several absolute accuracy thresholds without relabelling a lower-bound deviation diagnostic as a pass at a uniform upper-bound target. The experiments identify representation and search costs; they do not support an unqualified claim that the explicit actor is universally superior.

\subsection{B4: Continuous-to-discrete error}
We have removed the finite-action restriction from the new preference verification problem: the complete three-dimensional continuous control box is enclosed at every stored state and time. The verifier handles interpolation cells, reflected preference increments, binding action constraints, and wealth-exit switches. Unresolved action boxes retain their outward upper bounds. This replaces a missing action-coverage term by an actually computed bound for the nodal economy.

State and time transfer are separate. A new proposition states the optimal-operator and implemented-policy defects needed to compare the nodal economy with the continuously monitored economic diffusion, and proves their finite-horizon accumulation. We execute three nested space--time evaluations of each of six frozen policies with the same original time schedule and state-interpolated action rule. These are common-policy diagnostics, not an inferred error rate. The retained action-domain sensitivity study also distinguishes within-box refinement from changing the constrained economy. Continuous-state, diffusion-timing, and boundary-monitoring defects are not declared zero. Consequently neither a finite-game best response nor the new continuous-action nodal envelope is called a complete continuous-economy guarantee.

\subsection{B5: High-dimensional coverage and method pairing}
We recompute all 36 NBO-minus-direct comparisons from the pinned common-noise arrays: two dimensions, three seeds, two training domains, and three Euler resolutions. Each uses pathwise differences and a pointwise Student interval. Both narrow-box failures and wide-box improvements remain visible. We verify the stored common-noise identities rather than compare unrelated marginal confidence intervals.

The article continues to describe this as policy-evaluation and coverage evidence. Zero observed exits in a finite sample is not a zero exit probability; a Monte Carlo interval is not an optimal-value or time-discretization bound. The experiment preserves the unbounded capital economy instead of quietly changing it into a stopped training box.

\subsection{B6: Arithmetic replay and multidimensional verification}
The scalar result is described precisely as an independent interval-arithmetic replay on the previously generated complete cover. It reuses the stored weights and cover, and is not presented as independent training or a formal machine proof. Both the NBO and direct-HJB results and their costs are reported.

The new preference verifier extends the operational action account to all three continuous controls on a two-state nodal economy. Two enclosing implementations use cell-slope and cell-corner bounds, respectively. Both retain all unresolved upper bounds. A common optimal envelope can be reused for several frozen policies while each policy is independently evaluated. This is a concrete multidimensional action verification result, not a claim that unrestricted high-dimensional continuous-state verification has become inexpensive.

\subsection{B7: Coherence without changing the economic topic}
Rather than discard the economic agenda, we organize the paper around one chain: feasible policy proposal, independent policy evaluation, economically appropriate improvement or deviation, and a stated verification domain. The nonlinear consumption problem supplies the continuous differential certificate; endogenous preferences supplies continuous-action training and a complete action account; dynamic competition supplies full unilateral deviations; the coupled capital study supplies an explicit coverage diagnostic. Recursive utility, temporal selves, viscosity selection, and trace estimation remain connected to the same formulation. Exact-structure checks are clearly separated from neural training. The result is an integrated NBO revision, not a new paper obtained by deleting the difficult applications.

\section{Major Comments}
\subsection{M1: Name and training oracle}
The main text defines the evaluation--improvement map and positive-weight Bellman operator separately. It identifies the R7 all-action-supervised actor exactly, then specifies the distinct R8 continuous-output training route. The title remains \emph{Neural Bellman Operators}; no maximization advantage is inferred from the word ``operator.''

\subsection{M2: Raw, shortlist, and hybrid policies}
The R7 raw actor is the primary neural representation result. Raw, shortlist, and corrected policies have separate arrays, full dynamic checks, and correction frequencies. The R8 continuous raw actor and locally improved policy are likewise separate. Stored-table deployment, raw-network deployment, and online local improvement are not conflated. The supplement reports exact array bytes and a minimal correction-index payload, explicitly excluding container overhead and off-grid correctness.

\subsection{M3: Cost accounting}
New records separate construction where applicable, action optimization, local improvement, critic fitting, reference construction, policy evaluation, and full verification. All action queries, critic gradients, and L-BFGS closures are retained. Shared verification is charged explicitly. The R7 records do not contain every historical timing subdivision, so we do not invent one retrospectively. Cross-environment timing comparisons are labelled; equality of nominal update counts does not establish equal hardware cost or equal achieved accuracy.

\subsection{M4: Missing pure equilibria}
We independently locate all five missing pure fitted stages across the six R7 runs and print their state, time, and minimum maximum defect. They remain in the full dynamic verification. The supplement proves that pure Markov best-response maximization bounds randomized and history-dependent own deviations against fixed Markov rivals. This explains the deviation class actually checked without claiming that a mixed equilibrium profile was computed.

\subsection{M5: Paired payoff statistics}
Every stored dimension, seed, domain, and Euler resolution receives a paired mean, standard error, and pointwise 95 percent interval. The complete table has 36 rows. Shared paths across resolutions are acknowledged, and no simultaneous-coverage or continuous-time-bias conclusion is attached to these intervals.

\subsection{M6: Economic scale}
Utility and profit targets are stated in absolute units. The finite NDU value range and center magnitude are provided. The ratio of $.1$ to the center's absolute utility is a scale comparison, not a consumption-equivalent welfare percentage. Endogenous curvature and preference-adjustment costs preclude interpreting that ratio as a common consumption rescaling without a separate calculation. Game payoffs and the two market normalizations are identified explicitly.

\subsection{M7: Interval terminology}
We use the requested precise description of the independent scalar replay and retain the explicit distinction between outward arithmetic under stated assumptions and formal machine proof. The new action cover has its own scope, inputs, budgets, and retained unresolved boxes.

\subsection{M8: Development versus evidence}
The unexecuted post-evidence R7 development protocol is excluded as a source of numerical claims. New comparator configurations are disclosed as post-R7-review R8 development and actually executed. Initial smoke, raw-actor, multi-head, and direct-gradient failures remain alongside improved runs. The later tighter cover is an openly disclosed verification refinement, not retrospectively labelled an independent efficacy trial.

\subsection{M9: Authoritative candidate}
The final-candidate branch contains the integrated root manuscript and supplement, this point-by-point response, exact source and evidence identities, generated-table hashes, numerical and regression outputs, and compiled artifacts. Earlier review branches and historical paper files are not overwritten. The candidate's source and execution records identify what was run; a successful numerical or compilation check is not represented as referee acceptance.
